\section{Theoretical Results}
\label{sec:theory}

This section proves the Stage~1 results of
Section~\ref{subsec:linear_case}:
Lemmas~\ref{lem:foundation_regular_loss} and~\ref{lem:foundation_cover}
and Theorems~\ref{thm:foundation_generalization}
and~\ref{thm:foundation_verification}. The arguments are elementary and
are written to be checked line by line. Each proof opens with a short
statement of its idea. Sections~\ref{subsec:stage2_proofs}
and~\ref{subsec:stage3_envelope} add draft proofs for Stage~2 and a draft
Stage~3 starting point, both awaiting author review.

The results fit together as follows. Lemma~\ref{lem:foundation_regular_loss}
says that a small change of the certificate matrix makes only a small change
in the loss. Lemma~\ref{lem:foundation_cover} says that finitely many
matrices approximate every candidate. Theorem~\ref{thm:foundation_generalization}
combines the two with Hoeffding's inequality: concentration is proved at the
finitely many approximating matrices, and Lipschitz continuity carries it to
the whole class. Theorem~\ref{thm:foundation_verification} is independent of
the other three. It involves no data or probability and is the discrete-time
quadratic Lyapunov argument of course Chapter~3.

\subsection{Preliminaries}
\label{subsec:theory_prelims}

Let $\mathbb S^n$ denote the real symmetric $n\times n$ matrices, with
Frobenius inner product $\langle A,B\rangle_F=\operatorname{tr}(A^\top B)$
and norm $\|A\|_F=\langle A,A\rangle_F^{1/2}$. We use four standard facts.

\begin{enumerate}
\item[(F1)] \emph{Quadratic forms.} For $\Delta\in\mathbb S^n$ and
$y\in\mathbb R^n$,
\begin{equation}
\label{eq:theory_quadratic_form}
 |y^\top\Delta y|\le\|\Delta\|_F\|y\|_2^2 .
\end{equation}
Indeed, $y^\top\Delta y=\langle\Delta,yy^\top\rangle_F$, and Cauchy--Schwarz
gives $|\langle\Delta,yy^\top\rangle_F|\le\|\Delta\|_F\|yy^\top\|_F$, where
$\|yy^\top\|_F^2=\sum_{i,j}y_i^2y_j^2=\|y\|_2^4$.

\item[(F2)] \emph{Rayleigh quotient.} For $S\in\mathbb S^n$ the spectral
theorem gives an orthonormal eigenbasis, so
\begin{equation}
\label{eq:theory_rayleigh}
 x^\top Sx\ge\lambda_{\min}(S)\|x\|_2^2\quad\text{for all }x,
 \qquad\text{with equality at a unit eigenvector } v \text{ of }\lambda_{\min}(S).
\end{equation}
Applying the same fact to $-S$ gives $x^\top Sx\le\lambda_{\max}(S)\|x\|_2^2$.
Together these are the bound
$\lambda_{\min}(P)\|x\|^2\le x^\top Px\le\lambda_{\max}(P)\|x\|^2$ used in
course Chapter~3, Example~3.12.

\item[(F3)] \emph{Candidate bounds.} If $P\in\mathcal P$, then
$a\|x\|_2^2\le V_P(x)\le b\|x\|_2^2$ for every $x$. Hence every
eigenvalue of $P$ lies in $[a,b]$ (apply this at a unit eigenvector),
$\|P\|_2\le b$, and
\begin{equation}
\label{eq:theory_frobenius_bound}
 \|P\|_F^2=\operatorname{tr}(P^2)=\sum_{i=1}^n\lambda_i(P)^2\le nb^2 .
\end{equation}

\item[(F4)] \emph{Trajectories.} Within a trajectory the plant is fixed.
Hence $x_{t+1}=G_K(\theta)x_t$ and,
by~\eqref{eq:foundation_dynamics_bound},
$\|x_{t+1}\|_2\le M_K\|x_t\|_2$. In particular $x_t=0$ implies $x_s=0$ for all
$s\ge t$. The nonzero states among $x_0,\ldots,x_{H-1}$ therefore form an
initial segment $x_0,\ldots,x_{T-1}$, where
$T=T(\tau)=\#\{t<H:x_t\ne0\}\in\{0,\ldots,H\}$.
\end{enumerate}

\paragraph{Measurability.}
Give $\Theta$ its discrete $\sigma$-algebra and $\mathcal X_0$ its Borel
$\sigma$-algebra. Fix $P$. The map $z\mapsto s_P(\tau_K(z))$, with values
in $[-\infty,\infty)$, is measurable. Hence the loss $\ell_P(\tau_K(z))$
and the violation event $\{s_P(\tau_K(z))>0\}$ are measurable too, so
$L_K(P)$ and $e_{K,H}(P)$ are well defined.

To see this, first fix $\theta$ and write $G=G_K(\theta)$. The count
$T=\sum_{t<H}\mathbf 1\{G^tx_0\ne0\}$ is a sum of indicators of Borel sets,
so it is a measurable function of $x_0$. On $\{T=0\}$, $s_P=-\infty$. On
$\{T=k\}$ with $k\ge1$, $s_P$ is the maximum of the $k$ functions
$x_0\mapsto r_P(G^tx_0,G^{t+1}x_0)$, $t<k$. Each is continuous on the open
set $\{G^tx_0\ne0\}$, which contains $\{T=k\}$. So $s_P$ is Borel on each
of finitely many Borel pieces, and therefore Borel. Since $\Theta$ is
finite, the preimage of any Borel set under $z\mapsto s_P(\tau_K(z))$ is a
finite union of measurable rectangles $\{\theta\}\times(\text{Borel set})$,
which is measurable. Finally, $\ell_P=\phi(s_P)$ with the continuous $\phi$
of Step~3 below.

\subsection{Proof of Lemma~\ref{lem:foundation_regular_loss}}
\label{subsec:proof_regular_loss}

\emph{Idea.} The residual is linear in $P$, so changing $P$ changes it by a
quadratic form in $P-P'$, which (F1) controls. Taking a maximum and then
clipping are both Lipschitz operations, so the bound passes to the loss.

\begin{proof}
\emph{Step 1: residual.} Let $\Delta=P-P'\in\mathbb S^n$, let $x\ne0$, and
let $x^+=G_K(\theta)x$. The constant $\eta$ cancels in the difference:
\[
 r_P(x,x^+)-r_{P'}(x,x^+)
 =\frac{x^{+\top}\Delta x^+-x^\top\Delta x}{\|x\|_2^2}.
\]
By (F1) and $\|x^+\|_2\le M_K\|x\|_2$
(from~\eqref{eq:foundation_dynamics_bound}, since $\theta\in\Theta$), the
absolute value of the numerator is at most
\[
 \|\Delta\|_F(\|x^+\|_2^2+\|x\|_2^2)\le\|\Delta\|_F(M_K^2+1)\|x\|_2^2 .
\]
Dividing by $\|x\|_2^2$ gives~\eqref{eq:foundation_residual_lipschitz}.

\emph{Step 2: worst transition.} Fix a trajectory $\tau$. The set of times
$\{t<H:x_t\ne0\}$ is determined by $\tau$ alone and is the same for $P$ and
$P'$. If it is empty, then $\ell_P(\tau)=\ell_{P'}(\tau)=0$ by definition.
Otherwise $s_P(\tau)$ and $s_{P'}(\tau)$ are maxima over the same finite
index set. For real numbers, $|\max_t\alpha_t-\max_t\beta_t|\le
\max_t|\alpha_t-\beta_t|$. To see this, let $t^*$ attain $\max_t\alpha_t$;
then $\max_t\alpha_t-\max_t\beta_t\le\alpha_{t^*}-\beta_{t^*}$, and the other
direction is symmetric. Step~1 then gives
\[
 |s_P(\tau)-s_{P'}(\tau)|\le(M_K^2+1)\|\Delta\|_F .
\]

\emph{Step 3: clipping.} Write $\ell_P(\tau)=\phi(s_P(\tau))$, where
$\phi(s)=\min\{1,\max\{0,1+s/\gamma\}\}$ for real $s$ and
$\phi(-\infty)=0$. This matches the all-zero convention of
Definition~\ref{def:foundation_risk}. Step~2 has already handled that
case, so only real arguments matter below. The map $s\mapsto1+s/\gamma$ is
$(1/\gamma)$-Lipschitz, and clipping to $[0,1]$ is $1$-Lipschitz, so $\phi$
is $(1/\gamma)$-Lipschitz. Combining this with Step~2
gives~\eqref{eq:foundation_loss_lipschitz} with
$L_{\mathrm{par}}=(M_K^2+1)/\gamma$.

\emph{Step 4: range and domination.} The bound $0\le\ell_P\le1$ holds
because $\phi$ takes values in $[0,1]$. If $s_P(\tau)>0$, then
$1+s_P(\tau)/\gamma>1$, so $\ell_P(\tau)=1$. Otherwise the indicator is $0$,
which is at most $\ell_P(\tau)$. This covers the all-zero case
$s_P=-\infty$. So $\mathbf 1\{s_P(\tau)>0\}\le\ell_P(\tau)$ for every
$\tau$, and taking expectations under $\mathcal D_{\mathrm{cert}}$ gives
$e_{K,H}(P)\le L_K(P)$.
\end{proof}

Step~1 is the place where linear dynamics and the normalization by
$\|x\|_2^2$ are used. Because the bound is stated relative to $\|x\|_2^2$,
it does not depend on the size of the state or on the horizon $H$. This is
why Theorem~\ref{thm:foundation_generalization} has constants independent
of $H$. Step~1 holds for any symmetric $P,P'$, not only for members of
$\mathcal P$.

\subsection{Proof of Lemma~\ref{lem:foundation_cover}}
\label{subsec:proof_cover}

\emph{Idea.} We prove a volume (packing) bound in the $d$-dimensional
space $\mathbb S^n$. Course Proposition~5.18 states the analogous bound for
a Euclidean ball without proof. Course Proposition~5.20 rescales a cover of
the unit ball. Its centers are guaranteed to lie in the parameter set only
if that set contains the full ball of radius $B_\Theta$ about $0$.
Otherwise the rescaled points still cover the set, but they need not be an
$\varepsilon$-net in the sense of course Definition~5.15, which requires
centers in the set. Our $\mathcal P$ excludes a neighborhood of $0$, so it
does not contain such a ball. We therefore build the net from a maximal
separated subset of $\mathcal P$, which guarantees centers in $\mathcal P$.

\begin{proof}
\emph{Step 1: norm bound.} The first claim
in~\eqref{eq:foundation_cover} is (F3),
equation~\eqref{eq:theory_frobenius_bound}. Write $R=b\sqrt n$, so
$\mathcal P$ lies in the Frobenius ball of radius $R$ about $0$.

\emph{Step 2: coordinates.} The matrices $E_{ii}$ ($1\le i\le n$) and
$(E_{ij}+E_{ji})/\sqrt2$ ($i<j$), where $E_{ij}$ are the standard basis
matrices, form an orthonormal basis of $(\mathbb S^n,\langle\cdot,\cdot
\rangle_F)$. There are $n+n(n-1)/2=d$ of them. Expanding in this basis
identifies $\mathbb S^n$ with $\mathbb R^d$ and takes $\|\cdot\|_F$ to the
Euclidean norm. Let $\mathrm{vol}$ denote Lebesgue measure on $\mathbb R^d$.
A ball of radius $\rho$ then has volume $\rho^d\,\mathrm{vol}(B_1)$, where
$B_1$ is the unit ball.

\emph{Step 3: separated sets are small.} Call a finite set
$\{P_1,\ldots,P_k\}\subseteq\mathcal P$ \emph{$\varepsilon$-separated} if
$\|P_i-P_j\|_F>\varepsilon$ for all $i\ne j$. The open balls of radius
$\varepsilon/2$ about the $P_i$ are pairwise disjoint: a common point $Q$
would give $\|P_i-P_j\|_F\le\|P_i-Q\|_F+\|Q-P_j\|_F<\varepsilon$. Each of
these balls lies in the ball of radius $R+\varepsilon/2$ about $0$.
Comparing volumes,
\[
 k\Big(\frac\varepsilon2\Big)^d\mathrm{vol}(B_1)
 \le\Big(R+\frac\varepsilon2\Big)^d\mathrm{vol}(B_1),
 \qquad\text{so}\qquad
 k\le\Big(1+\frac{2R}{\varepsilon}\Big)^d .
\]

\emph{Step 4: a maximal separated set is a net.} The set $\mathcal P$ is
nonempty, since $aI\in\mathcal P$. By Step~3 the sizes of
$\varepsilon$-separated subsets are bounded, so we may choose one of largest
size, say $\{P_1,\ldots,P_k\}$. Every $P\in\mathcal P$ satisfies
$\|P-P_i\|_F\le\varepsilon$ for some $i$. Otherwise $P$ could be added to
the set and it would stay $\varepsilon$-separated, contradicting
maximality. Hence $\{P_i\}$ is a Frobenius $\varepsilon$-net with centers in
$\mathcal P$, and $\mathcal N_F(\varepsilon,\mathcal P)\le k\le
(1+2b\sqrt n/\varepsilon)^d$. Taking logarithms completes the proof.
\end{proof}

Steps~3--4, with $\mathcal P$ replaced by the unit ball of $\mathbb R^n$,
give course Proposition~5.18. The dimension that matters is that of the
symmetric matrices, $d=n(n+1)/2$. This matches the parameter count
$k=d(d+1)/2$ of the class (4.3) of \citet{boffi2021learning}, whose $d$ is
the feature dimension; with the feature map $m(x)=x$ their feature
dimension is our $n$, and their $k$ is our $d$. The proof
of Theorem~\ref{thm:foundation_generalization} would also work with centers
outside $\mathcal P$, since the loss is defined and bounded for every
symmetric matrix. Keeping the centers in $\mathcal P$ makes each center an
admissible certificate.

\subsection{Proof of Theorem~\ref{thm:foundation_generalization}}
\label{subsec:proof_generalization}

\emph{Idea.} This follows course Theorem~5.8 (uniform convergence over a
finite class), applied to the finite net of Lemma~\ref{lem:foundation_cover}.
Lemma~\ref{lem:foundation_regular_loss} then extends the result from the net
to all of $\mathcal P$, at a cost of $L_{\mathrm{par}}\varepsilon$ on each
of the two sides. As in course Section~5.6.2, one sample is one whole
trajectory, so the sample size is $N$ rather than $NH$.

\begin{proof}
Write $\Delta_N=\Delta_N(\varepsilon,\delta)$ and
\[
 u=\sqrt{\frac{d\log(1+2b\sqrt n/\varepsilon)+\log(2/\delta)}{2N}}>0,
 \qquad\text{so that}\qquad \Delta_N=2L_{\mathrm{par}}\varepsilon+u .
\]

\emph{Step 1: a fixed net.} Using Lemma~\ref{lem:foundation_cover}, choose
an $\varepsilon$-net $\{P^1,\ldots,P^{M_\varepsilon}\}\subseteq\mathcal P$
with $M_\varepsilon\le(1+2b\sqrt n/\varepsilon)^d$. The net depends only on $\mathcal P$ and
$\varepsilon$, which are fixed before sampling, so it does not depend on the
data. This is what allows Hoeffding's inequality to be applied at each
center (course Remark~5.7).

\emph{Step 2: concentration at one center.} Fix $j$ and set
$X_i=\ell_{P^j}(\tau_i)$, for $i=1,\ldots,N$. Each $X_i$ is a fixed
measurable function of $z_i$. The $z_i$ are i.i.d.\ draws from
$\mathcal D_{\mathrm{cert}}$, so the $X_i$ are independent, lie in $[0,1]$
by Lemma~\ref{lem:foundation_regular_loss}, and have mean $L_K(P^j)$. Their
average is $\widehat L_{N,K}(P^j)$. Course Lemma~5.6, applied once to the
$X_i$ and once to the $-X_i$, gives
\[
 \Pr\big\{|\widehat L_{N,K}(P^j)-L_K(P^j)|\ge u\big\}
 \le2\exp(-2Nu^2).
\]

\emph{Step 3: union bound.} Let $E$ be the event that
$|\widehat L_{N,K}(P^j)-L_K(P^j)|<u$ for every $j=1,\ldots,M_\varepsilon$.
As in the proof of course Theorem~5.8, but with a two-sided bound (hence
the factor $2$),
\[
 \Pr(E^c)\le2M_\varepsilon\exp(-2Nu^2)
 =2M_\varepsilon\cdot\frac{\delta}{2}
   \Big(1+\frac{2b\sqrt n}{\varepsilon}\Big)^{-d}
 \le\delta .
\]
So $\Pr(E)\ge1-\delta$. The event $E$ involves only finitely many
measurable functions of the sample, so it is measurable.

\emph{Step 4: from the net to all of $\mathcal P$.} Assume $E$ occurs, and
let $P\in\mathcal P$ be arbitrary. Choose $j$ with
$\|P-P^j\|_F\le\varepsilon$. By Lemma~\ref{lem:foundation_regular_loss},
$|\ell_P(\tau)-\ell_{P^j}(\tau)|\le L_{\mathrm{par}}\varepsilon$ for every
trajectory $\tau$. This is the pointwise form of the Lipschitz transfer in
course Proposition~5.19. Averaging the bound over the sample, or taking its
expectation, gives
\[
 |\widehat L_{N,K}(P)-\widehat L_{N,K}(P^j)|\le L_{\mathrm{par}}\varepsilon,
 \qquad
 |L_K(P)-L_K(P^j)|\le L_{\mathrm{par}}\varepsilon .
\]
The triangle inequality and the definition of $E$ then give
\[
 |L_K(P)-\widehat L_{N,K}(P)|
 \le L_{\mathrm{par}}\varepsilon+u+L_{\mathrm{par}}\varepsilon=\Delta_N .
\]
Since $P$ was arbitrary and $E$ does not depend on $P$, this
proves~\eqref{eq:foundation_uniform_bound} simultaneously for all
$P\in\mathcal P$.

\emph{Step 5: the fitted certificate.} On $E$, the bound of Step~4 holds for
every member of $\mathcal P$, including any $\widehat P$ chosen from the
data. Lemma~\ref{lem:foundation_regular_loss} gives
$e_{K,H}(\widehat P)\le L_K(\widehat P)\le\widehat L_{N,K}(\widehat P)
+\Delta_N$, and $e_{K,H}(\widehat P)\le1$ because it is a probability. This
is~\eqref{eq:foundation_fitted_risk}.

\emph{Step 6: near-optimality.} Assume also~\eqref{eq:foundation_erm}. For
any $P\in\mathcal P$, on $E$,
\begin{align*}
 L_K(\widehat P)
 &\le\widehat L_{N,K}(\widehat P)+\Delta_N
 \le\inf_{P'\in\mathcal P}\widehat L_{N,K}(P')+\xi_{\mathrm{opt}}+\Delta_N\\
 &\le\widehat L_{N,K}(P)+\xi_{\mathrm{opt}}+\Delta_N
 \le L_K(P)+2\Delta_N+\xi_{\mathrm{opt}} .
\end{align*}
The first and last inequalities use Step~4. Taking the infimum over $P$
gives~\eqref{eq:foundation_erm_population}.

\emph{Step 7: rate.} Put $\varepsilon=N^{-1/2}$ and $c=2b\sqrt n$, and let
$N\ge1$. Since $1+c\sqrt N\le(1+c)(1+N)$ and $\log(1+N)\ge\log2$,
\[
 \log(1+c\sqrt N)\le\log(1+c)+\log(1+N)
 \le\Big(1+\frac{\log(1+c)}{\log2}\Big)\log(1+N).
\]
Write $\kappa=\log(1+c)/\log2$ and
$Q_N=d\log(1+N)+\log(1/\delta)$. Since $d\ge1$,
$\log(2/\delta)\le\log(1+N)+\log(1/\delta)\le Q_N$. The numerator of $u$
is therefore at most $d(1+\kappa)\log(1+N)+Q_N\le(2+\kappa)Q_N$, so
$u\le\sqrt{1+\kappa/2}\,\sqrt{Q_N/N}$. For the first term of $\Delta_N$,
using $\log(1+N)\ge\log2$, $d\ge1$, and $\log(1/\delta)>0$,
\[
 2L_{\mathrm{par}}N^{-1/2}
 \le\frac{2L_{\mathrm{par}}}{\sqrt{\log2}}\sqrt{\frac{d\log(1+N)}{N}}
 \le\frac{2L_{\mathrm{par}}}{\sqrt{\log2}}\sqrt{\frac{Q_N}{N}} .
\]
Adding the two bounds, for every $N\ge1$ and $\delta\in(0,1)$,
\[
 \Delta_N\le C\sqrt{\frac{d\log(1+N)+\log(1/\delta)}{N}},
 \qquad
 C=\frac{2L_{\mathrm{par}}}{\sqrt{\log2}}
   +\sqrt{1+\frac{\log(1+2b\sqrt n)}{2\log2}} .
\]
The constant $C$ depends only on $b$, $n$, and $L_{\mathrm{par}}$.
\end{proof}

Two features of the proof deserve comment. First, we read ``with
probability at least $1-\delta$'' as: there is a measurable event $E$ with
$\Pr(E)\ge1-\delta$ on which all the stated bounds hold. The theorem's
phrase ``on the same event'' uses this reading. Step~4 never takes a
supremum over the uncountable class $\mathcal P$, so no measurability of a
supremum is needed, and~\eqref{eq:foundation_fitted_risk}
and~\eqref{eq:foundation_erm_population} do not need a measurable fitting
map. (The supremum in~\eqref{eq:foundation_uniform_bound} is in fact
measurable: both losses are $L_{\mathrm{par}}$-Lipschitz in $P$, so it
equals a supremum over a countable dense subset of $\mathcal P$.) The
measurability required in Definition~\ref{def:foundation_fitting} is needed
only for statements about random quantities such as
$\mathbb E[e_{K,H}(\widehat P)]$.

Second, the proof uses only that the $N$ trajectories are independent with
law $\mathcal D_{\mathrm{cert}}$, and that the net and the constants
$K,\varepsilon,\delta,\gamma,\eta$ are fixed in advance. A data-dependent
$\varepsilon$ breaks Step~1, because the net would depend on the data. A
data-dependent $K$, $\gamma$, $\eta$, or sampling law breaks Step~2,
because the $X_i$ would no longer be fixed functions of i.i.d.\ draws with
mean $L_K(P^j)$. This is why the theorem must be re-applied with fresh data
after the Stage~2 update.

\subsection{Proof of Theorem~\ref{thm:foundation_verification}}
\label{subsec:proof_verification}

\emph{Idea.} For a linear system, the change in $V_P$ along one step is the
quadratic form of $-(P-G^\top PG)$, by course Example~3.12,
equation~(3.22). That example has $Q\succ0$ from a Lyapunov equation; here
the matrix $P-G^\top PG$ plays the role of $Q$, and it is not assumed
positive definite in advance. By (F2), the smallest eigenvalue of
$P-G^\top PG$ is the exact worst-case normalized decrease. A uniform decrease, together with the
sandwich $a\|x\|^2\le V_P\le b\|x\|^2$, makes $V_P$ contract geometrically.
This is course Lemma~3.11 with $p=2$.

\begin{proof}
Fix a retained plant $\theta$. Write $G=G_K(\theta)$ and
$S_\theta=P-G^\top PG\in\mathbb S^n$, so that
$\eta_\theta(P)=\lambda_{\min}(S_\theta)$. For every $x$,
\begin{equation}
\label{eq:theory_drift_identity}
 V_P(Gx)-V_P(x)=x^\top G^\top PGx-x^\top Px=-x^\top S_\theta x .
\end{equation}

\emph{Step 1: the eigenvalue test.} By~\eqref{eq:theory_drift_identity},
condition~\eqref{eq:foundation_desired_decrease} holds for all $x$ if and
only if $x^\top S_\theta x\ge\eta\|x\|_2^2$ for all $x$. If
$\lambda_{\min}(S_\theta)\ge\eta$, this follows from (F2). Conversely, at a
unit eigenvector $v$ of $\lambda_{\min}(S_\theta)$, the inequality reads
$\lambda_{\min}(S_\theta)=v^\top S_\theta v\ge\eta$.

\emph{Step 2: $0<\eta_*\le b$.} Positivity is assumed. The minimum
in~\eqref{eq:foundation_verified_margin} is attained, since
$\Theta_{\mathrm{ret}}$ is finite and nonempty. Let $u$ be a unit
eigenvector of $\lambda_{\min}(P)$. For any retained $\theta$, (F2) and
$P\succeq0$ give
\[
 \eta_*\le\eta_\theta(P)\le u^\top S_\theta u
 =u^\top Pu-(Gu)^\top P(Gu)\le u^\top Pu=\lambda_{\min}(P)\le b .
\]
The sharper bound $\eta_*\le\lambda_{\min}(P)$ shows that the case
$\eta_*=b$ is degenerate. It forces $\lambda_{\min}(P)=b$, hence $P=bI$
because all eigenvalues lie in $[a,b]$. Then
$S_\theta=b(I-G^\top G)$ has smallest eigenvalue $b(1-\|G\|_2^2)$, which
equals $b$ only if $G_K(\theta)=0$ for every retained plant.

\emph{Step 3: geometric contraction.} For every retained $\theta$ and every
$x$, equation~\eqref{eq:theory_drift_identity} and (F2) give
$V_P(Gx)\le V_P(x)-\eta_\theta\|x\|_2^2\le V_P(x)-\eta_*\|x\|_2^2$. By (F3),
$\|x\|_2^2\ge V_P(x)/b$, so
\[
 V_P(Gx)\le\Big(1-\frac{\eta_*}{b}\Big)V_P(x),
 \qquad 0\le1-\frac{\eta_*}{b}<1,
\]
where the range of the factor comes from Step~2. Write
$\varrho=1-\eta_*/b$. Induction on $t$ gives
$V_P(x_t)\le\varrho^{\,t}V_P(x_0)$. Applying (F3) on both sides,
\[
 a\|x_t\|_2^2\le V_P(x_t)\le\varrho^{\,t}V_P(x_0)\le\varrho^{\,t}b\|x_0\|_2^2,
\]
which is~\eqref{eq:foundation_exponential_bound}. The bound also holds at
$t=0$, with the convention $\varrho^{\,0}=1$, because $b/a>1$. Taking
square roots,
$\|x_t\|_2\le\sqrt{b/a}\,(\sqrt\varrho)^{t}\|x_0\|_2$ for every initial
state. This is global exponential stability in the sense of course
Definition~3.8, with $C=\sqrt{b/a}$ and rate $\sqrt\varrho\in[0,1)$. In the
notation of course Lemma~3.11 this is the case $p=2$,
$\bar\alpha_1=a$, $\bar\alpha_2=b$, $\bar\alpha_3=\eta_*$. That lemma
needs $\bar\alpha_3>0$ (so that $\bar\alpha_3r^2$ is a class~$\mathcal K$
function), $\bar\alpha_1\le\bar\alpha_2$, and $\bar\alpha_3\le\bar\alpha_2$.
These are the hypothesis $\eta_*>0$, the assumption $a<b$, and Step~2. If $\eta_*=b$, then $\varrho=0$, so
$V_P(x_1)\le0$, and $x_1=0$ because $V_P(x_1)\ge a\|x_1\|_2^2$.

\emph{Step 4: invariant ellipsoids.} If $V_P(x)\le c$, then
$V_P(Gx)\le V_P(x)-\eta_*\|x\|_2^2\le c$. So each sublevel set
$\mathcal E_c(P)$ is mapped into itself by every retained closed loop.

\emph{Step 5: zero violation risk.} If $\eta_*\ge\eta$, then
$\eta_\theta(P)\ge\eta$ for every retained $\theta$. By Step~1, the desired
decrease holds at every state, so $r_P(x,G_K(\theta)x)\le0$ for every
$x\ne0$. Any trajectory of a retained plant therefore has $s_P(\tau)\le0$,
or $s_P(\tau)=-\infty$ if all its states are zero. Because $\Theta$ is
discrete, $\Theta_{\mathrm{ret}}\times\mathcal X_0$ is closed, so
``supported on'' means
$\mathcal D_{\mathrm{cert}}(\Theta_{\mathrm{ret}}\times\mathcal X_0)=1$. Then
$e_{K,H}(P)\le\Pr\{\theta\notin\Theta_{\mathrm{ret}}\}=0$, and this holds for
every horizon $H$.
\end{proof}

The proof uses the known matrices $A(\theta)$ and $B(\theta)$, not
sampled data. Its conclusion is therefore deterministic and does not
depend on how $P$ was obtained, including fitting by
Definition~\ref{def:foundation_fitting}. Fixing $\Theta_{\mathrm{ret}}$
before inspecting results is a protocol requirement, not a mathematical
hypothesis: the proof never uses that timing. What it protects is the
honesty of the support condition in Step~5. Since
$\mathcal D_{\mathrm{cert}}$ is a law on $\Theta\times\mathcal X_0$, that
condition is the same as requiring the $\Theta$-marginal to put mass~$1$
on $\Theta_{\mathrm{ret}}$. The fact that the decrease holds at every
state, not only on $\mathcal X_0$, would matter only if the initial-state
law were later extended beyond $\mathcal X_0$.

In floating point the theorem is applied to computed quantities, and three
checks are needed. First, the tolerance must cover the error in forming
$S_\theta=P-G^\top PG$ as well as the eigensolver error; by Weyl's
inequality the computed eigenvalue moves by at most the norm of the error
in $S_\theta$. Second, the certified lower bound $\widehat\eta_\theta-
\mathrm{tol}$, not $\widehat\eta_\theta$ itself, should be used: it must
be positive for stability, at least $\eta$ for the claim of Step~5, and it
is the value to use in the rate $\varrho$. Third, a fitted $P$ must itself
be checked to be symmetric with $aI\preceq P\preceq bI$, since (F3) and
Step~2 assume $P\in\mathcal P$ exactly.

\subsection{Draft proofs of the Stage~2 statements}
\label{subsec:stage2_proofs}

These proofs are drafts awaiting author review. Throughout, $G=G_K(\theta)$,
$E=-B(\theta)\Delta K$, and $G^+=G+E=G_{K^+}(\theta)$.

\begin{proof}[Proof of Lemma~\ref{lem:stage2_margin_drift}]
Expanding $(G+E)^\top P(G+E)-P$ gives~\eqref{eq:stage2_drift_expansion}. Let
$T=G^\top PE+E^\top PG+E^\top PE$, which is symmetric. For a unit vector $y$,
\[
 y^\top Ty=2y^\top G^\top PEy+y^\top E^\top PEy
 \le2\|G^\top P\|_2\|E\|_2+\|P\|_2\|E\|_2^2=c_\theta(\Delta K),
\]
because $\|E\|_2=d_\theta(\Delta K)$. By (F2), $T\preceq c_\theta(\Delta K)I$,
which is~\eqref{eq:stage2_drift_bound}. Since
$G^\top PG-P\preceq-\eta_\theta I$, an admissible change gives
$(G^+)^\top PG^+-P\preceq-(\eta_\theta-c_\theta(\Delta K))I
\preceq-\rho\eta_\theta I$.
\end{proof}

\begin{proof}[Proof of Theorem~\ref{thm:stage2_certified_descent}]
\emph{Accepted steps.} An accepted nonzero step has $D\ne0$, $s_j>0$, and
$\Delta K=s_jD$ admissible. Armijo~\eqref{eq:stage2_armijo} with $c_A>0$ gives
$F_q(K^+)<F_q(K)$. Lemma~\ref{lem:stage2_margin_drift} gives
\eqref{eq:stage2_preserved_margin}. Write
$\eta_*^+=\min_{\theta\in\Theta_{\mathrm{ret}}}
\lambda_{\min}(P-G_{K^+}(\theta)^\top PG_{K^+}(\theta))$ and
$\underline\eta=\rho\min_{\theta\in\Theta_{\mathrm{ret}}}\eta_\theta$, so
$\eta_*^+\ge\underline\eta>0$. Theorem~\ref{thm:foundation_verification}, applied to
$K^+$ with the same $P\in\mathcal P$, gives exponential stability, invariance of
every $\mathcal E_c(P)$, and~\eqref{eq:foundation_exponential_bound} with
$\eta_*^+$. Since $0<\underline\eta\le\eta_*^+\le b$, replacing $\eta_*^+$ by
$\underline\eta$ only enlarges the right side. If $\rho\eta_\theta\ge\eta$ on the
retention set, its eigenvalue test retains the coefficient $\eta$.

\emph{Small steps.} Let $D\ne0$ and $\varphi(s)=F_q(K+sD)$. Each $G_{K+sD}(\theta)$
is affine in $s$, so $\varphi$ is a polynomial, and
$\varphi'(0)=\langle\nabla_KF_q(K),D\rangle_F=-\|D\|_F^2$. Thus
$(\varphi(s)-\varphi(0))/s\to-\|D\|_F^2<-c_A\|D\|_F^2$ as $s\downarrow0$, since
$c_A<1$, so Armijo holds on some $(0,s_A]$. The radius holds for
$s\le\kappa/\|D\|_2$. Each $s\mapsto c_\theta(sD)$ is continuous and vanishes at
$0$, each budget $(1-\rho)\eta_\theta$ is positive, and $\Theta_{\mathrm{ret}}$ is
finite, so the consumption test holds on some $(0,s_G]$. As $s_j\to0$, an uncapped
search reaches a passing trial unless it has already accepted one. The fallback
$K^+=K$ changes nothing.
\end{proof}

\emph{Starting from the step capacity.} The midterm report and the linear studies
start the search at $\bar s=\min\{s_0,s^*(D)\}$, where $s^*(D)$ is the step capacity of
Proposition~\ref{prop:stage3_direction}, and shrink only for Armijo. The accepted-step
argument above is unchanged, because every trial is admissible and passes Armijo. The
small-step argument is also unchanged: $\bar s>0$, all smaller steps are admissible, and
Armijo holds on $(0,s_A]$. The implementation multiplies the first trial by
$1-10^{-9}$ so that rounding keeps it inside the guard.

\begin{proof}[Proof of Lemma~\ref{lem:stage2_cost_sensitivity}]
Let $\|\Delta K\|_2\le\kappa$ and $W=W_K$, $W^+=W_{K^+}$. Then
$\|G\|_2,\|G^+\|_2\le\overline M_\theta$ and $\|E\|_2\le\|B(\theta)\|_2\|\Delta K\|_2$.
Split
\begin{align*}
 P_H(K^+,\theta)-P_H(K,\theta)
 &=\sum_{t=0}^{H-1}\big((G^+)^t\big)^\top(W^+-W)(G^+)^t\\
 &\quad+\sum_{t=1}^{H-1}\Big(\big((G^+)^t\big)^\top W(G^+)^t-(G^t)^\top WG^t\Big).
\end{align*}
For the first sum, $W^+-W=K^\top R\Delta K+\Delta K^\top RK+\Delta K^\top R\Delta K$
gives $\|W^+-W\|_2\le\|R\|_2(2\|K\|_2+\kappa)\|\Delta K\|_2$. For the second, write
$A=(G^+)^t$ and $C=G^t$. Then
\[
 A^\top WA-C^\top WC=A^\top W(A-C)+(A-C)^\top WC,
 \qquad
 A-C=\sum_{j=0}^{t-1}(G^+)^{t-1-j}EG^j,
\]
so $\|A-C\|_2\le t\overline M_\theta^{\,t-1}\|E\|_2$, and the $t$-th term has norm at
most $2\|W\|_2\|B(\theta)\|_2t\overline M_\theta^{\,2t-1}\|\Delta K\|_2$. Summing
gives~\eqref{eq:stage2_cost_matrix_drift}. Then
$|x^\top(P_H(K^+,\theta)-P_H(K,\theta))x|\le r^2\mathcal L_H(\theta;\kappa)
\|\Delta K\|_2$ on $\mathcal X_0$, which is~\eqref{eq:stage2_cost_drift}.
\end{proof}

\begin{proof}[Proof of Theorem~\ref{thm:stage2_cost_regression}]
Write $b=b_J(\theta;\Delta K)$. If $J_{H,K}(\theta,x)\le h-b$,
then~\eqref{eq:stage2_cost_drift} gives $J_{H,K^+}(\theta,x)\le h$,
proving~\eqref{eq:stage2_interior_preservation}. If $z$ is lost, then
$J_{H,K}(z)\le h<J_{H,K^+}(z)\le J_{H,K}(z)+b$, so $z$ lies in the loss band,
proving~\eqref{eq:stage2_losses}. Costs are continuous in $x$ and $\Theta$ is finite,
so these sets are measurable. Finally
$\mu(S_{K^+}^J)=\mu(S_{K^+}^J\cap S_K^J)+\mu(S_{K^+}^J\setminus S_K^J)$, likewise
for $S_K^J$; subtracting gives the identity
in~\eqref{eq:stage2_coverage_accounting}, and dropping the gained mass
with~\eqref{eq:stage2_losses} gives the inequality.
\end{proof}

\begin{proof}[Proof of Proposition~\ref{prop:stage2_counterexample}]
Here $G_{\widetilde K}(\theta)=\theta-\widetilde K$ and
$J_{2,\widetilde K}(\theta,1)=1+(\theta-\widetilde K)^2$. At $K=1/2$ both costs are
$5/4=h$, so $C_J(K)=1$. With $q=(1/4,3/4)$,
$F_q(\widetilde K)=\tfrac14(1+\widetilde K^2)+\tfrac34(1+(1-\widetilde K)^2)$, so
$F_q(1/2)=1.25$, $\nabla F_q(1/2)=1/4-3/4=-1/2$, and $D=1/2$. Then $s_0=2/5$ gives
$\Delta K=1/5\le\kappa$ and $K^+=7/10$. Both old margins are $1-1/4=3/4$. With
$d_\theta=1/5$, $c_\theta=2(1/2)(1/5)+1/25=0.24\le0.375=(1-\rho)(3/4)$.
Also $F_q(7/10)=\tfrac14(1.49)+\tfrac34(1.09)=1.19\le1.25-\tfrac12\cdot\tfrac25
\cdot\tfrac14=1.20$, so the first trial passes. The new costs are $1.49>h$ at
$\theta=0$ and $1.09\le h$ at $\theta=1$, so $C_J(K^+)=1/2$. The new margins are
$1-0.49=0.51$ and $1-0.09=0.91$.
\end{proof}

\subsection{Draft Stage~3 starting point: envelope accounting}
\label{subsec:stage3_envelope}

This proposition is a draft awaiting author review. Fix $K$, a direction $D\ne0$, and
$\bar s=\kappa/\|D\|_2$. Keep $\overline M_\theta$ from
Lemma~\ref{lem:stage2_cost_sensitivity}, and set $\beta_\theta=\|B(\theta)D\|_2$,
$\bar k=\|K\|_2+\kappa$, $w_0=\|Q\|_2+\|R\|_2\bar k^2$,
$w_1=2\|R\|_2\bar k\|D\|_2$, and $w_2=2\|R\|_2\|D\|_2^2$. For $z=(\theta,x)$, let
$X_0=\|x\|_2$, $U_0=V_0=0$, and
\[
 X_{t+1}=\overline M_\theta X_t,\quad
 U_{t+1}=\overline M_\theta U_t+\beta_\theta X_t,\quad
 V_{t+1}=\overline M_\theta V_t+2\beta_\theta U_t,
\]
\[
 C_z=\sum_{t=0}^{H-1}\big(2w_0(U_t^2+X_tV_t)+4w_1X_tU_t+w_2X_t^2\big),
\]
and define the \emph{alignment} and \emph{upper envelope}
\[
 a_z=-\langle\nabla_KJ_{H,K}(z),D\rangle_F,
 \qquad \mathcal U_z(s)=J_{H,K}(z)-sa_z+\tfrac12C_zs^2.
\]

\begin{proposition}[Envelope coverage bound]
\label{prop:stage3_envelope}
For every $s\in[0,\bar s]$, $J_{H,K+sD}(z)\le\mathcal U_z(s)$. Define the
\emph{guaranteed gains} and \emph{possible losses}
\[
 \mathcal A(s)=\{z\notin S_K^J:\mathcal U_z(s)\le h\},
 \qquad
 \mathcal R(s)=\{z\in S_K^J:\mathcal U_z(s)>h\}.
\]
Then
\[
 C_J(K+sD)-C_J(K)\ge\mu(\mathcal A(s))-\mu(\mathcal R(s)).
\]
Suppose $C_z>0$. A rejected $z$ with excess $e_z=J_{H,K}(z)-h>0$ is in
$\mathcal A(s)$ if and only if $a_z>0$, $a_z^2\ge2C_ze_z$, and
$s\in[(a_z-\sqrt{a_z^2-2C_ze_z})/C_z,(a_z+\sqrt{a_z^2-2C_ze_z})/C_z]$. An accepted
$z$ with slack $\sigma_z=h-J_{H,K}(z)\ge0$ is outside $\mathcal R(s)$ if and only if
$s\le(a_z+\sqrt{a_z^2+2C_z\sigma_z})/C_z$.
\end{proposition}

\begin{proof}
Let $r\in[0,\bar s]$, $K_r=K+rD$, and $G_r=G_K(\theta)-rB(\theta)D$, so
$\|G_r\|_2\le\overline M_\theta$ and $\|K_r\|_2\le\bar k$. Let $x_t=G_r^tx$, with
derivatives in $r$ marked by primes. Then
$x_{t+1}'=G_rx_t'-B(\theta)Dx_t$ and
$x_{t+1}''=G_rx_t''-2B(\theta)Dx_t'$, so induction gives $\|x_t\|_2\le X_t$,
$\|x_t'\|_2\le U_t$, and $\|x_t''\|_2\le V_t$. For $W_r=Q+K_r^\top RK_r$ we have
$\|W_r\|_2\le w_0$, $\|W_r'\|_2\le w_1$, and $\|W_r''\|_2\le w_2$. Since
\[
 (x_t^\top W_rx_t)''=2x_t'^\top W_rx_t'+2x_t^\top W_rx_t''
 +4x_t^\top W_r'x_t'+x_t^\top W_r''x_t,
\]
$\phi(r)=J_{H,K_r}(z)$ satisfies $|\phi''(r)|\le C_z$. Also $\phi'(0)=-a_z$, and
Taylor's formula with integral remainder gives
$\phi(s)\le\mathcal U_z(s)$. Thus $\mathcal A(s)$ is contained in the gained set, and
the lost set is contained in $\mathcal R(s)$. The bound follows
from~\eqref{eq:stage2_coverage_accounting}'s identity. The last two claims solve
$\tfrac12C_zs^2-a_zs+e_z\le0$ and $\tfrac12C_zs^2-a_zs-\sigma_z\le0$ over $s\ge0$.
For the first, the roots have positive product, so real roots are both positive
exactly when $a_z>0$.
\end{proof}

This bound uses old costs and gradients only, so it is computable before updated
costs are observed. Positivity is not automatic. On the 201 primary-family updates of
the recorded study it is never positive (largest value $-0.156$), although coverage
rises by up to $0.203$ on those updates. On the designed one-plant scalar fixture,
with $A=0.8$, $B=1$, $K=0.2$, $Q=1$, $R=0$, $H=2$, $h=1.2$, and four resets, the
boundary and high-cost targets give a bound of $0.25$, equal to the observed change,
while background gives $0$. This is a step-size effect, not evidence of target
superiority: with one plant every cost gradient has the same sign, and background's
shorter move $\Delta K\approx0.14$ leaves the largest cost at $1.21$, just above $h$,
whereas boundary's $\Delta K\approx0.22$ crosses it.

\begin{proposition}[Guard-limited updates depend only on direction]
\label{prop:stage3_direction}
Fix $K$ and a certificate $P$ with positive retained margins $\eta_\theta$, and let
\[
 \mathcal G=\{\Delta K:\|\Delta K\|_2\le\kappa,\
 c_\theta(\Delta K)\le(1-\rho)\eta_\theta
 \text{ for every }\theta\in\Theta_{\mathrm{ret}}\}.
\]
For $D\ne0$, let $s^*(D)=\max\{s\ge0:sD\in\mathcal G\}$, the \emph{step capacity}.
Then $s^*(D)$ is finite and positive, $\{s\ge0:sD\in\mathcal G\}=[0,s^*(D)]$, and for
every $\alpha>0$,
\[
 s^*(\alpha D)\,\alpha D=s^*(D)\,D.
\]
So the full-capacity update depends on $D$ only through $D/\|D\|_F$.
\end{proposition}

\begin{proof}
$\mathcal G$ is defined by conditions on $\Delta K$ alone. Along $s\mapsto sD$, the
norm $s\|D\|_2$ and each $c_\theta(sD)=2\|G_K(\theta)^\top P\|_2sd_\theta(D)
+\|P\|_2s^2d_\theta(D)^2$ are continuous and nondecreasing in $s\ge0$. They vanish at
$s=0$, where every budget is positive, and $s\|D\|_2$ exceeds $\kappa$ for large $s$.
Hence the feasible set is a compact interval $[0,s^*(D)]$ with $s^*(D)>0$. Finally,
$s\alpha D\in\mathcal G$ exactly when $s\alpha\le s^*(D)$, so
$s^*(\alpha D)=s^*(D)/\alpha$.
\end{proof}

The proposition ignores two limits on the step: the initial step $s_0$ and the Armijo
test. When either binds, the size of $D$ can matter. In the recorded family study the
$s_0$ limit almost never binds, but Armijo limits most updates on mirrored plants.
Under the proposition's conditions, a target can change the update only by turning
$D$, so any theorem that separates target rules needs a hypothesis on that turn.
